\documentclass[12pt]{article}
\usepackage{amsmath,amssymb,amsthm}

\begin{document}
	\subsection*{Formal Definitions}
	
	\paragraph{Symmetric Group. $(S_n, \circ) $}  
	Let $X$ be a finite set with $n$ elements, for example $X = \{1,2,\dots,n\}$.  
	The \emph{symmetric group on $n$ letters}, denoted $S_n$, is defined as
	$$
	S_n = \{ \sigma : X \to X \;\mid\; \sigma \text{ is a bijection} \}.
	$$
	
	The group operation is composition of functions:
	$$
	(\sigma \circ \tau)(x) = \sigma(\tau(x)), \quad x \in X.
	$$
	
	The identity element is the identity map $\mathrm{id}_X$, and the inverse of $\sigma$ is the inverse bijection $\sigma^{-1}$.  
	The cardinality is
	$$
	|S_n| = n!.
	$$
	
	---
	
	\paragraph{Automorphism Group of $S_n$.}  
	For a group $G$, an \emph{automorphism} is a bijective group homomorphism
	$$
	\varphi : G \to G, \qquad \varphi(xy) = \varphi(x)\varphi(y).
	$$
    
    \subsection*{Two-Level Perspective on $S_n$}
    
    There are two levels involved when working with the symmetric group:
    
    \begin{itemize}
    	\item \textbf{Level 1 (Set level):}  
    	The underlying set is 
    	$$
    	X = \{1,2,\dots,n\}.
    	$$
    	At this level, a permutation is a bijection 
    	$$
    	\sigma : X \to X.
    	$$
    	
    	\item \textbf{Level 2 (Group level):}  
    	The elements of the group $S_n$ are the permutations themselves.  
    	Thus
    	$$
    	S_n = \{\sigma : X \to X \;\mid\; \sigma \text{ is bijective}\}.
    	$$
    	The group operation is \emph{composition of permutations}:
    	$$
    	(\sigma \circ \tau)(x) = \sigma(\tau(x)), \qquad \sigma,\tau \in S_n.
    	$$
    \end{itemize}
    
    \noindent
    Hence:
    \begin{itemize}
    	\item The \emph{basic elements} of $S_n$ are the permutations, not the numbers $1,\dots,n$.  
    	\item The group operation is not a simple swap of numbers, but a \emph{composition of bijections}, producing new permutations.
    \end{itemize}
    
	
	The \emph{automorphism group of $S_n$} is defined as
	$$
	\mathrm{Aut}(S_n) = \{ \varphi : S_n \to S_n \;\mid\; \varphi \text{ is a bijective homomorphism} \}.
	$$
	
	Elements of $\mathrm{Aut}(S_n)$ are bijections of $S_n$ onto itself that preserve the group law.  
	Composition of automorphisms is the group operation.
	
	\medskip
	
	\noindent \textbf{Key fact:} For $n \neq 6$,
	$$
	\mathrm{Aut}(S_n) \cong S_n,
	$$
	since every automorphism is given by conjugation inside $S_n$.  
	For $n=6$, $\mathrm{Aut}(S_6)$ is strictly larger due to the existence of outer automorphisms.
	
	\section*{Symmetric Groups and Automorphisms}
	
	\subsection*{Case $S_1$}
	The symmetric group $S_n$ is defined as the group of all bijections (permutations) of the set
	$$
	\{1,2,\dots,n\}.
	$$
	
	For $S_1$, the underlying set is $\{1\}$.  
	There is only one bijection:
	$$
	\sigma : 1 \mapsto 1,
	$$
	so
	$$
	S_1 = \{\mathrm{id}_{\{1\}}\}.
	$$
	This is the trivial group.
	
	---
	
	\subsection*{Permutations vs Automorphisms of $S_n$}
	- Each $\sigma \in S_n$ is a bijection of the underlying set $\{1,\dots,n\}$.  
	- This makes it an \emph{automorphism of the set}.  
	- But an \emph{automorphism of the group $S_n$} means a bijective group homomorphism
	$$
	\varphi : S_n \to S_n, \qquad \varphi(\sigma\tau) = \varphi(\sigma)\varphi(\tau).
	$$
	
	So:
	- $S_n$ = permutations of the set of $n$ elements.  
	- $\mathrm{Aut}(S_n)$ = automorphisms of the group $S_n$ itself.  
	
	---
	
	\subsection*{Inner Automorphisms}
	Most automorphisms of $S_n$ are given by conjugation:
	$$
	\varphi_g(\sigma) = g\sigma g^{-1}, \qquad g \in S_n.
	$$
	
	- For $n \neq 6$: all automorphisms are inner, so $\mathrm{Aut}(S_n) \cong S_n$.  
	- For $n=6$: there exist outer automorphisms, making $\mathrm{Aut}(S_6)$ larger.
	
	---
	
	\subsection*{Homomorphisms, Isomorphisms, Automorphisms}
	- A group homomorphism: $\varphi: G \to H$, $\;\varphi(xy) = \varphi(x)\varphi(y)$.  
	- An isomorphism: a bijective homomorphism.  
	- An endomorphism: a homomorphism $G \to G$.  
	- An automorphism: a bijective endomorphism.  
	
	Thus $\mathrm{Aut}(G)$ is a group under composition.
	
	---
	
	\subsection*{Case $S_2$}
	Underlying set: $\{1,2\}$.  
	Two bijections:
	- Identity: $e: 1\mapsto 1, 2\mapsto 2$.  
	- Transposition: $\tau: 1\mapsto 2, 2\mapsto 1$.
	
	So
	$$
	S_2 = \{ e, (12)\}.
	$$
	
	Multiplication table:
	$$
	\begin{array}{c|cc}
		\circ & e & (12) \\ \hline
		e & e & (12) \\
		(12) & (12) & e
	\end{array}
	$$
	
	Thus $S_2 \cong C_2$.  
	Since $C_2$ is abelian of order 2, $\mathrm{Aut}(S_2)$ is trivial.
	
	---
	
	\subsection*{Interpretation: Labels vs Slots}
	- As swapping \emph{labels}: $(12)$ relabels $1 \leftrightarrow 2$.  
	- As swapping \emph{slots}: $(12)$ moves the element in slot 1 to slot 2 and vice versa.  
	
	Both interpretations represent the same permutation.
	
	---
	
	\subsection*{Case $S_3$}
	Underlying set: $\{1,2,3\}$.  
	
	Elements:
	$$
	S_3 = \{ (1)(2)(3),\; (12),\; (13),\; (23),\; (123),\; (132)\}.
	$$
	
	- Identity: $(1)(2)(3)$.  
	- 3 transpositions: $(12),(13),(23)$.  
	- 2 three-cycles: $(123),(132)$.  
	
	Non-abelian: for example,
	$$
	(12)(23) = (123), \qquad (23)(12) = (132).
	$$
	
	Subgroups:
	- $\langle (123)\rangle \cong C_3$.  
	- Each transposition generates a $C_2$.  
	
	Automorphism group: $\mathrm{Aut}(S_3) \cong S_3$.
	
	---
	
	\subsection*{Formal Definition of a Cycle}
	Let $X=\{1,\dots,n\}$ and $\sigma \in S_n$.  
	An $r$-cycle is a permutation
	$$
	\sigma = (a_1\; a_2\; \dots\; a_r),
	$$
	with distinct $a_i \in X$, such that
	$$
	\sigma(a_1)=a_2,\; \sigma(a_2)=a_3,\;\dots,\;\sigma(a_r)=a_1,
	$$
	and $\sigma(x)=x$ for all $x \notin \{a_1,\dots,a_r\}$.
	
	Every permutation can be uniquely written (up to order) as a product of disjoint cycles.
	
	---
	
	\subsection*{Counting $r$-cycles}
	Number of $r$-cycles in $S_n$:
	$$
	\binom{n}{r}(r-1)! = \frac{1}{r} \, {}_nP_r = \frac{n!}{r(n-r)!}.
	$$
	
	Reason:  
	- Choose $r$ elements: $\binom{n}{r}$.  
	- Arrange them in a circle: $(r-1)!$ ways.  
	- Equivalent to $\frac{1}{r} {}_nP_r$ since cycles are invariant under rotation.
	
	---
	
	\subsection*{Check in $S_3$}
	- $r=1$: identity element.  
	- $r=2$: $\frac{1}{2} {}_3P_2 = 3$ transpositions: $(12),(13),(23)$.  
	- $r=3$: $\frac{1}{3} {}_3P_3 = 2$ three-cycles: $(123),(132)$.  
	
	Total: $1+3+2 = 6$ elements.
	
	---
	
	\subsection*{Cycles as Circular Permutations}
	- Linear permutations: ${}_nP_r = \dfrac{n!}{(n-r)!}$.  
	- Circular permutations: divide by $r$ (rotations are the same cycle).  
	
	Thus cycles in $S_n$ are exactly circular permutations of $r$ chosen elements.
	
	\subsection*{Observation on Cycle Notation}
	
	In cycle notation, the ordering of elements is relevant only up to rotation or reversal:
	
	\begin{itemize}
		\item A cycle written with different starting points represents the \textbf{same} permutation:
		$$
		(1\,2\,3) = (2\,3\,1) = (3\,1\,2).
		$$
		\item Writing a cycle backwards gives its \textbf{inverse}:
		$$
		(1\,2\,3)^{-1} = (1\,3\,2).
		$$
		\item Disjoint cycles commute, so their order in a product is irrelevant:
		$$
		(12)(34) = (34)(12).
		$$
	\end{itemize}
	
	Therefore, when counting cycles, what matters is capturing all distinct \emph{combinations} of elements, without overcounting mere repetitions due to rotation. This is why the number of $r$-cycles in $S_n$ is given by
	$$
	\frac{1}{r}\, {}_nP_r = \frac{n!}{r (n-r)!}.
	$$
	
	
	---
	
	\subsection*{Geometric Model: $S_3$ as Symmetries of a Triangle}
	Take an equilateral triangle with vertices labeled $1,2,3$.  
	
	\begin{itemize}
		\item \textbf{Rotations}:
		\begin{align*}
			& (1)(2)(3) \quad (\text{identity}), \\
			& (123) \quad (\text{rotation by } 120^\circ), \\
			& (132) \quad (\text{rotation by } 240^\circ).
		\end{align*}
		\item \textbf{Reflections}:
		\begin{align*}
			& (12), \\
			& (13), \\
			& (23).
		\end{align*}
	\end{itemize}
	
	So
	$$
	\mathrm{Isom}(\triangle) \;\cong\; S_3 \;\cong\; D_6.
	$$
	
	---
	
	\subsection*{Careful about Geometry}
	Only \emph{isometries} (rigid motions) are allowed:
	- rotations and reflections.  
	- permutations that deform the triangle are not valid isometries.  
	
	For $n=3$: every permutation is an isometry, so $S_3 \cong D_6$.  
	For $n \geq 4$: only $2n$ of the $n!$ permutations correspond to isometries.
	
	---
	
	\subsection*{General Case}
	- $S_n$: all permutations of $n$ labels ($n!$ elements).  
	- $D_{2n}$: isometry group of a regular $n$-gon ($2n$ elements).  
	
	Thus:
	$$
	S_3 \cong D_6, \qquad
	D_{2n} \subset S_n \;\text{ for } n \geq 4.
	$$
	
	---
	
	\subsection*{Convention on Composition in $S_n$}
	
	By definition of function composition,
	$$
	(\sigma \circ \tau)(x) = \sigma(\tau(x)),
	$$
	so the rightmost permutation acts first.
	
	\medskip
	
	\noindent\textbf{Example in $S_5$:}
	\[
	\sigma = (1\,5\,2\,4\,3) = (15)(52)(24)(43).
	\]
	
	We compute step by step, applying the rightmost transposition first:
	
	\begin{itemize}
		\item Start with $x=1$.
		\begin{align*}
			(43): &\; 1 \mapsto 1, \\
			(24): &\; 1 \mapsto 1, \\
			(52): &\; 1 \mapsto 1, \\
			(15): &\; 1 \mapsto 5.
		\end{align*}
		So overall $1 \mapsto 5$.
		
		\item Start with $x=5$.
		\begin{align*}
			(43): &\; 5 \mapsto 5, \\
			(24): &\; 5 \mapsto 5, \\
			(52): &\; 5 \mapsto 2, \\
			(15): &\; 2 \mapsto 2.
		\end{align*}
		So overall $5 \mapsto 2$.
		
		\item Start with $x=2$.
		\begin{align*}
			(43): &\; 2 \mapsto 2, \\
			(24): &\; 2 \mapsto 4, \\
			(52): &\; 4 \mapsto 4, \\
			(15): &\; 4 \mapsto 4.
		\end{align*}
		So overall $2 \mapsto 4$.
		
		\item Start with $x=4$.
		\begin{align*}
			(43): &\; 4 \mapsto 3, \\
			(24): &\; 3 \mapsto 3, \\
			(52): &\; 3 \mapsto 3, \\
			(15): &\; 3 \mapsto 3.
		\end{align*}
		So overall $4 \mapsto 3$.
		
		\item Start with $x=3$.
		\begin{align*}
			(43): &\; 3 \mapsto 4, \\
			(24): &\; 4 \mapsto 2, \\
			(52): &\; 2 \mapsto 5, \\
			(15): &\; 5 \mapsto 1.
		\end{align*}
		So overall $3 \mapsto 1$.
	\end{itemize}
	
	Putting this together, we recover the 5-cycle:
	$$
	1 \mapsto 5 \mapsto 2 \mapsto 4 \mapsto 3 \mapsto 1,
	$$
	that is,
	$$
	(15)(52)(24)(43) = (1\,5\,2\,4\,3).
	$$
	
		
		\section*{Analysis of the Symmetric Group $S_3$}
		
		\subsection*{1. Underlying Set}
		
		Let 
		\[
		X = \{1,2,3\}.
		\]
		The symmetric group on $X$ is defined as
		\[
		S_3 = Aut_{set}(X),
		\]
		where $Aut_{set}(X)$ denotes the set of all bijections $f: X \to X$.  
		This is equivalent to $\mathrm{Sym}(X)$, the full permutation group of $X$.
		
		\subsection*{2. Group Operation}
		
		For $f,g \in S_3$, the group operation is composition of functions:
		\[
		(g \cdot f)(x) = g(f(x)), \qquad \forall x \in X.
		\]
		Thus, $(S_3,\cdot)$ is a group with identity the identity map $id_X$ and inverses given by $f^{-1}$ for each $f \in S_3$.
		
		\subsection*{3. Action of $f$ on $X$}
		
		Consider $f \in Aut_{set}(X)$ defined by
		\[
		f(1)=2, \quad f(2)=1, \quad f(3)=3.
		\]
		This bijection swaps $1$ and $2$ while fixing $3$.  
		Applied to the ordered triple $(1,2,3)$, it produces
		\[
		(1,2,3) \mapsto (2,1,3).
		\]
		
		\subsection*{4. Action of $g$ on $X$}
		
		Now consider $g \in Aut_{set}(X)$ defined by
		\[
		g(1)=3, \quad g(2)=1, \quad g(3)=2.
		\]
		This is a $3$-cycle, cycling the elements as $1 \mapsto 3 \mapsto 2 \mapsto 1$.  
		Applied to $(1,2,3)$, it produces
		\[
		(1,2,3) \mapsto (3,1,2).
		\]
		
		\subsection*{5. Composition $g \cdot f$}
		
		To compute the composition $g \cdot f$, apply $f$ first and then $g$:
		
		\[
		\begin{aligned}
			x=1: &\quad f(1)=2, \quad g(f(1))=g(2)=1, \\
			x=2: &\quad f(2)=1, \quad g(f(2))=g(1)=3, \\
			x=3: &\quad f(3)=3, \quad g(f(3))=g(3)=2.
		\end{aligned}
		\]
		
		Thus,
		\[
		(g \cdot f)(1)=1, \quad (g \cdot f)(2)=3, \quad (g \cdot f)(3)=2.
		\]
		
		In tuple form:
		\[
		(g \cdot f) : (1,2,3) \mapsto (1,3,2).
		\]
		
		\subsection*{6. Remarks on Arbitrariness}
		
		The choice of $f$ and $g$ is arbitrary: any bijection $X \to X$ is a valid element of $S_3$.  
		What is fixed is the rule of the group operation:
		\[
		(g \cdot f)(x) = g(f(x)).
		\]
		Once $f$ and $g$ are specified, their composition is uniquely determined.  
		
		The “looping” observed arises naturally from the cycle structure of permutations. For example, $g=(1\,3\,2)$ means following the images of $1$, $3$, and $2$ leads back to the starting point, closing a cycle.
		
	\end{document}
	
	
\end{document}
